\centeredsection{АННОТАЦИЯ}
Работа посвящена практической реализации компилятора базисного подмножества языка
Рефал-5 с особым вниманием к представлению объектных выражений в памяти. Семейство
языков, основанных на переписывании термов, традиционно сталкивается с двумя
конкурирующими проблемами: дорогим копированием значений переменных при списковом
представлении и дорогой конкатенацией при массивном. Каждая известная реализация
решает один из этих компромиссов за счёт другого, и попытка обойти оба сразу
оказывается нетривиальной инженерной задачей.

В качестве представления объектных выражений в работе выбрана раскрученная
двусторонняя очередь, описанная Окасаки. Эта структура амортизированно даёт
константную сложность как для операций на концах выражения, так и для конкатенации,
а структурное разделение между версиями устраняет физическое копирование при
многократном использовании переменной. Компилятор переводит исходный текст через
лексический, синтаксический и семантический анализ во внутреннее представление
программы, а исполняющий модуль вычисляет нормализованную форму, опираясь на
выбранную структуру данных. Реализация выполнена на языке Java~25~LTS, сборка
автоматизирована средствами Apache~Maven, корректность подтверждена набором
автоматических тестов JUnit~5 и сравнительным прогоном с двумя альтернативными
реализациями Рефала-5: классической Refal-5~PZ на языке~C и Refal-5J на JVM.

Результаты работы могут использоваться как самостоятельный компилятор базисного
подмножества Рефала-5 и как платформа для дальнейших исследований эффективных
представлений объектных выражений.


Работа состоит из 64~страниц, включает 4~раздела, 4~приложения, 9~таблиц,
4~рисунка и 36~листингов. Библиографический список содержит 13~источников.
